%=============================================================================
% MODULE 02: ABSTRACT (Final Expansion)
%=============================================================================
% Frontiers in Public Health — Structured abstract (250-400 words)
%=============================================================================

\begin{abstract}
\noindent
\textbf{Background:} Machine learning models are increasingly used for 
diabetes diagnosis, yet algorithmic bias threatens equitable healthcare 
delivery. Existing research has predominantly focused on single demographic 
attributes, leaving intersectional effects poorly understood.

\textbf{Objective:} This study evaluates algorithmic bias in diabetes 
prediction models across multiple demographic dimensions, examining 
performance disparities by gender, ethnicity, and their intersection.

\textbf{Methods:} We analyzed the Pima Indians Diabetes Database (n=768) 
using four machine learning algorithms: Logistic Regression, Random Forest, 
Gradient Boosting, and Support Vector Machine. Models were evaluated using 
5-fold stratified cross-validation and held-out test sets. Fairness was 
assessed using Equalized Odds Difference, Demographic Parity Difference, 
True Positive Rate Disparity, and False Positive Rate Disparity. 
Intersectional analysis examined performance at the intersection of gender 
and ethnicity.

\textbf{Results:} All four models exhibited measurable bias. The SVM model 
achieved the best balance between performance (F1=0.8313, AUC=ROC=0.8794) 
and fairness. False positive rates were consistently higher for minority 
patients (0.333--0.417) compared to majority patients (0.164--0.200). 
Intersectional analysis revealed amplified bias for minority females, who 
exhibited the lowest True Positive Rate (0.714) and F1-Score (0.750) among 
all subgroups. Cross-validation demonstrated stable performance across folds 
(SD < 0.04).

\textbf{Conclusions:} Algorithmic bias in diabetes prediction models is 
systemic and multidimensional. Single-attribute analyses underestimate 
the extent of discrimination. Intersectional approaches are essential for 
identifying compounded disadvantage. We recommend multi-faceted bias 
mitigation strategies spanning data collection, algorithm development, 
and post-deployment monitoring.

\noindent
\textbf{Keywords:} algorithmic bias, diabetes, machine learning, fairness, 
intersectionality, healthcare equity, machine learning ethics, bias 
mitigation
\end{abstract}
